% ============================================================
%  CORRIGÉ — FICHE DE RÉVISION — SÉRIES DE FOURIER — 2e BTS
%  Style : activite_fourier_a0 / Lycée Denis Diderot
% ============================================================

\documentclass[a4paper,11pt]{article}
\usepackage[utf8]{inputenc}
\usepackage[T1]{fontenc}
\usepackage[french]{babel}
\usepackage{fouriernc}
\usepackage{amsmath,amssymb,mathtools}
\usepackage{geometry}
\usepackage{tikz}
\usepackage{enumitem}
\usepackage{fancyhdr}
\usepackage[most]{tcolorbox}
\usepackage{pgfplots}
\pgfplotsset{compat=1.18}
\usepackage{xcolor}
\usepackage{array}
\usepackage{booktabs}
\usepackage{lastpage}

\geometry{margin=2cm}
\setlength{\headheight}{14pt}

\pagestyle{fancy}
\fancyhf{}
\lhead{BTS -- Mathématiques}
\rhead{Séries de Fourier -- \textcolor{red!70!black}{\textbf{CORRIGÉ}}}
\cfoot{\thepage}

\tcbset{
  exobox/.style={colback=orange!5!white, colframe=orange!60!black,
    fonttitle=\bfseries, title=#1, breakable},
  corrigebox/.style={colback=green!5!white, colframe=green!50!black,
    sharp corners, boxrule=0.5pt, left=6pt, right=6pt, top=4pt, bottom=4pt,
    breakable}
}

\newcommand{\R}{\mathbb{R}}
\newcommand{\Z}{\mathbb{Z}}
\newcommand{\w}{\omega_0}

% ==============================================================
\begin{document}
% ==============================================================

\begin{center}
  {\LARGE \textbf{Corrigé -- Séries de Fourier}} \\[0.2cm]
  {\large \textcolor{red!70!black}{\textbf{--- CORRIGÉ ---}}} \\[0.2cm]
  \rule{0.65\textwidth}{0.4pt}
\end{center}

\vspace{0.4cm}

% ==============================================================
\begin{tcolorbox}[exobox={Exercice 1 — Calcul des coefficients (cas général)}]
Fonction $f$ périodique de période $T = 2$ : $f(t) = 3$ sur $[0,1[$ et $f(t) = -1$ sur $[1,2[$.
\end{tcolorbox}

\begin{enumerate}[label=\textbf{\arabic*.}]

\item \textbf{Pulsation fondamentale}
\begin{tcolorbox}[corrigebox]
$\omega_0 = \dfrac{2\pi}{T} = \dfrac{2\pi}{2} = \pi\ \text{rad/s}$
\end{tcolorbox}

\item \textbf{Parité de $f$}
\begin{tcolorbox}[corrigebox]
Sur $[-1\,;\,1]$, pour $t\in\,]0,1[$: $f(-t) = f(-t+2) = f(2-t) = -1$ (car $2-t\in\,]1,2[$).\\
Or $f(t) = 3 \neq -1$, donc $f(-t)\neq f(t)$ et $f(-t) \neq -f(t)$.\\
La fonction $f$ est \textbf{ni paire ni impaire} : on doit calculer $a_0$, $a_n$ et $b_n$.
\end{tcolorbox}

\item \textbf{Calcul de $a_0$}
\begin{tcolorbox}[corrigebox]
\[
a_0 = \frac{1}{T}\int_0^T f(t)\,\mathrm{d}t
= \frac{1}{2}\left[\int_0^1 3\,\mathrm{d}t + \int_1^2 (-1)\,\mathrm{d}t\right]
= \frac{1}{2}\bigl[3\cdot 1 + (-1)\cdot 1\bigr] = \frac{1}{2}\times 2 = \boxed{1}
\]
$a_0 = 1$ est la \textbf{valeur moyenne} (composante continue) du signal.
\end{tcolorbox}

\item \textbf{Calcul de $a_n$}
\begin{tcolorbox}[corrigebox]
\begin{align*}
a_n &= \frac{2}{2}\left[\int_0^1 3\cos(n\pi t)\,\mathrm{d}t + \int_1^2 (-1)\cos(n\pi t)\,\mathrm{d}t\right] \\
&= 3\left[\frac{\sin(n\pi t)}{n\pi}\right]_0^1 - \left[\frac{\sin(n\pi t)}{n\pi}\right]_1^2
= \frac{3\underbrace{\sin(n\pi)}_{=0}}{n\pi} - \frac{\underbrace{\sin(2n\pi)}_{=0}-\underbrace{\sin(n\pi)}_{=0}}{n\pi} = 0
\end{align*}
$\boxed{a_n = 0}$ pour tout $n \geq 1$ (les bornes annulent tous les sinus).
\end{tcolorbox}

\item \textbf{Calcul de $b_n$}
\begin{tcolorbox}[corrigebox]
\begin{align*}
b_n &= \frac{2}{2}\left[\int_0^1 3\sin(n\pi t)\,\mathrm{d}t + \int_1^2 (-1)\sin(n\pi t)\,\mathrm{d}t\right] \\
&= 3\left[-\frac{\cos(n\pi t)}{n\pi}\right]_0^1 - \left[-\frac{\cos(n\pi t)}{n\pi}\right]_1^2 \\
&= \frac{3\bigl(1-\cos(n\pi)\bigr)}{n\pi} - \frac{\cos(n\pi)-\cos(2n\pi)}{n\pi} \\
&= \frac{3-3(-1)^n - (-1)^n + 1}{n\pi} = \boxed{\frac{4\bigl(1-(-1)^n\bigr)}{n\pi}}
\end{align*}
\begin{itemize}[nosep]
  \item $n$ \textbf{pair} : $b_n = 0$
  \item $n$ \textbf{impair} : $b_n = \dfrac{8}{n\pi}$
\end{itemize}
$b_1 = \dfrac{8}{\pi} \approx 2{,}55$ \qquad $b_2 = 0$ \qquad $b_3 = \dfrac{8}{3\pi} \approx 0{,}85$
\end{tcolorbox}

\item \textbf{Développement jusqu'à $n = 3$}
\begin{tcolorbox}[corrigebox]
\[
  S(t) = 1 + \frac{8}{\pi}\sin(\pi t) + \frac{8}{3\pi}\sin(3\pi t) + \cdots
\]
Seuls les harmoniques impairs sont présents ($a_n=0$ et $b_n=0$ pour $n$ pair).
\end{tcolorbox}

\end{enumerate}

\newpage

% ==============================================================
\begin{tcolorbox}[exobox={Exercice 2 — Parité et Parseval}]
Fonction $f$ \textbf{paire}, $T = 2\pi$, $\omega_0 = 1$ : $f(t) = 5$\,V sur $[0,\pi/2]$ et $f(t) = 0$ sur $]\pi/2,\pi]$.
\end{tcolorbox}

\begin{enumerate}[label=\textbf{\arabic*.}]

\item \textbf{Justification $b_n = 0$}
\begin{tcolorbox}[corrigebox]
$f$ est paire, $\sin(n\omega_0 t)$ est impaire. Le produit $f(t)\sin(n\omega_0 t)$ est donc \textbf{impair}.
Son intégrale sur une période symétrique $[-\pi,\pi]$ est nulle. Donc $b_n = 0$ pour tout $n \geq 1$.
\end{tcolorbox}

\item \textbf{Calcul de $a_0$}
\begin{tcolorbox}[corrigebox]
Pour $f$ paire : $a_0 = \dfrac{2}{T}\displaystyle\int_0^{T/2}f(t)\,\mathrm{d}t = \frac{2}{2\pi}\int_0^{\pi}f(t)\,\mathrm{d}t$
\[
a_0 = \frac{1}{\pi}\left[\int_0^{\pi/2} 5\,\mathrm{d}t + \int_{\pi/2}^{\pi} 0\,\mathrm{d}t\right]
= \frac{1}{\pi}\cdot\frac{5\pi}{2} = \boxed{\frac{5}{2} = 2{,}5\ \text{V}}
\]
$a_0 = 2{,}5$ V est la \textbf{valeur moyenne} du signal (composante continue).
\end{tcolorbox}

\item \textbf{Calcul de $a_n$}
\begin{tcolorbox}[corrigebox]
Avec $\omega_0 = 1$ (car $T = 2\pi$) :
\[
a_n = \frac{4}{2\pi}\int_0^{\pi}f(t)\cos(nt)\,\mathrm{d}t
= \frac{2}{\pi}\int_0^{\pi/2} 5\cos(nt)\,\mathrm{d}t
= \frac{10}{\pi}\left[\frac{\sin(nt)}{n}\right]_0^{\pi/2}
= \boxed{\frac{10}{n\pi}\sin\!\left(\frac{n\pi}{2}\right)}
\]
\[
a_1 = \frac{10}{\pi}\sin\!\left(\frac{\pi}{2}\right) = \frac{10}{\pi} \approx 3{,}18
\qquad
a_2 = \frac{10}{2\pi}\sin(\pi) = 0
\qquad
a_3 = \frac{10}{3\pi}\sin\!\left(\frac{3\pi}{2}\right) = -\frac{10}{3\pi} \approx -1{,}06
\]
\end{tcolorbox}

\item \textbf{Développement jusqu'à $n = 3$}
\begin{tcolorbox}[corrigebox]
\[
S(t) = \frac{5}{2} + \frac{10}{\pi}\cos(t) - \frac{10}{3\pi}\cos(3t) + \cdots
\]
\end{tcolorbox}

\item \textbf{Parseval}

\begin{enumerate}[label=\textbf{\alph*)}]
\item \textbf{$F_{\mathrm{eff}}$ par intégration directe}
\begin{tcolorbox}[corrigebox]
Sur $[0, 2\pi]$ : $f = 5$ sur $[0,\pi/2]$ et $[3\pi/2,2\pi]$ (par parité et periodicité), nulle ailleurs.
\[
F_{\mathrm{eff}}^2 = \frac{1}{2\pi}\left[25\cdot\frac{\pi}{2} + 0\cdot\pi + 25\cdot\frac{\pi}{2}\right]
= \frac{25\pi}{2\pi} = \frac{25}{2}
\qquad\Longrightarrow\qquad
\boxed{F_{\mathrm{eff}} = \frac{5}{\sqrt{2}} \approx 3{,}54\ \text{V}}
\]
\end{tcolorbox}

\item \textbf{Vérification par Parseval (tronqué à $n = 3$)}
\begin{tcolorbox}[corrigebox]
\[
F_{\mathrm{eff}}^2 \approx a_0^2 + \frac{1}{2}\left(a_1^2 + a_2^2 + a_3^2\right)
= \left(\frac{5}{2}\right)^{\!2} + \frac{1}{2}\left[\left(\frac{10}{\pi}\right)^{\!2} + 0 + \left(\frac{10}{3\pi}\right)^{\!2}\right]
\approx 6{,}25 + 5{,}63 \approx 11{,}88
\]
Valeur exacte : $12{,}5$. \quad Écart relatif : $\dfrac{12{,}5 - 11{,}88}{12{,}5} \approx 5\,\%$
\end{tcolorbox}
\end{enumerate}

\end{enumerate}

\newpage

% ==============================================================
\begin{tcolorbox}[exobox={Exercice 3 — Vrai ou Faux ?}]
\end{tcolorbox}

\renewcommand{\arraystretch}{2.2}
\begin{tabular}{|p{11.5cm}|>{\centering\arraybackslash}p{1.5cm}|>{\centering\arraybackslash}p{1.5cm}|}
\hline
\textbf{Affirmation} & \textbf{Vrai} & \textbf{Faux} \\
\hline
Si $f$ est \textbf{impaire}, alors $a_0 = 0$ et $a_n = 0$ pour tout $n \geq 1$. & \textcolor{green!60!black}{\textbf{X}} & \\
\hline
$a_0$ représente la valeur moyenne de $f$ sur une période. & \textcolor{green!60!black}{\textbf{X}} & \\
\hline
Une fonction de période $T = 0{,}02$\,s a pour fréquence $f_0 = 50$\,Hz. & \textcolor{green!60!black}{\textbf{X}} & \\
\hline
Le spectre en amplitude d'un signal périodique est un spectre \textbf{continu}. & & \textcolor{red!70!black}{\textbf{X}} \\
\hline
Si $f$ vérifie $f(x+T/2) = -f(x)$ pour tout $x$, alors tous les harmoniques de rang \textbf{pair} sont nuls. & \textcolor{green!60!black}{\textbf{X}} & \\
\hline
\end{tabular}

\vspace{0.3cm}

% ==============================================================
\begin{tcolorbox}[exobox={Exercice 4 — Spectre et Parseval}]
Signal créneau impair, amplitude $E$, période $T$ :
$f(t) = \dfrac{4E}{\pi}\displaystyle\sum_{\substack{n=1\\n\,\text{impair}}}^{+\infty}\dfrac{1}{n}\sin(n\omega_0 t)$
\end{tcolorbox}

\begin{enumerate}[label=\textbf{\arabic*.}]

\item \textbf{Harmoniques présents}
\begin{tcolorbox}[corrigebox]
Seuls les harmoniques de rang \textbf{impair} ($n = 1, 3, 5, 7, \ldots$) sont présents ($b_n \neq 0$).\\
Le créneau vérifie $f(t+T/2) = -f(t)$ (symétrie demi-onde), ce qui annule tous les coefficients de rang pair.
\end{tcolorbox}

\item \textbf{Amplitudes et spectre} ($E=1$, $T=1$ s, $\omega_0 = 2\pi$ rad/s)
\begin{tcolorbox}[corrigebox]
$a_n = 0$ ; $b_n = \dfrac{4}{n\pi}$ pour $n$ impair, donc $C_n = \sqrt{a_n^2+b_n^2} = \dfrac{4}{n\pi}$.\\[4pt]
$C_1 = \dfrac{4}{\pi} \approx 1{,}27$ \qquad $C_3 = \dfrac{4}{3\pi} \approx 0{,}42$ \qquad $C_5 = \dfrac{4}{5\pi} \approx 0{,}25$ \qquad $C_7 \approx 0{,}18$\\[4pt]
$C_0 = C_2 = C_4 = C_6 = 0$ \quad (rang pair ou nul).
\end{tcolorbox}

\begin{center}
\begin{tikzpicture}[x=1.5cm,y=3cm]
  \draw[->] (-0.3,0) -- (8,0) node[right]{\small$n$};
  \draw[->] (0,-0.05) -- (0,1.55) node[above]{\small$C_n$};
  \foreach \n in {0,1,...,7} \draw (\n,2pt)--(\n,-2pt) node[below]{\small$\n$};
  \foreach \y/\lab in {0.42/{0{,}42},0.85/{0{,}85},1.27/{1{,}27}}
    \draw (2pt,\y)--(-2pt,\y) node[left]{\tiny$\lab$};
  \draw[gray!25,xstep=1,ystep=0.42] (0,0) grid (7.5,1.4);
  \fill[orange!40] (0.7,0) rectangle (1.3,1.273);
  \fill[orange!40] (2.7,0) rectangle (3.3,0.424);
  \fill[orange!40] (4.7,0) rectangle (5.3,0.255);
  \fill[orange!40] (6.7,0) rectangle (7.3,0.182);
  \draw[orange!80!black,thick] (0.7,0) rectangle (1.3,1.273);
  \draw[orange!80!black,thick] (2.7,0) rectangle (3.3,0.424);
  \draw[orange!80!black,thick] (4.7,0) rectangle (5.3,0.255);
  \draw[orange!80!black,thick] (6.7,0) rectangle (7.3,0.182);
  \node[above,font=\tiny,orange!80!black] at (1,1.273) {$1{,}27$};
  \node[above,font=\tiny,orange!80!black] at (3,0.424) {$0{,}42$};
  \node[above,font=\tiny,orange!80!black] at (5,0.255) {$0{,}25$};
  \node[above,font=\tiny,orange!80!black] at (7,0.182) {$0{,}18$};
  \foreach \n in {0,2,4,6}
    \draw[blue!50!black,thick] (\n-0.3,0)--(\n+0.3,0);
\end{tikzpicture}
\end{center}

\item \textbf{Démonstration de $\displaystyle\sum_{\substack{n\geq 1\\n\,\text{impair}}}\dfrac{1}{n^2} = \dfrac{\pi^2}{8}$}
\begin{tcolorbox}[corrigebox]
\textbf{Calcul direct de $F_{\mathrm{eff}}$} (avec $E=1$, $T=1$\,s ; $f = \pm 1$ p.p.) :
\[
F_{\mathrm{eff}}^2 = \frac{1}{1}\int_0^1 1^2\,\mathrm{d}t = 1
\]
\textbf{Parseval} ($a_n=0$ et $b_n = \dfrac{4}{n\pi}$ pour $n$ impair, nul sinon) :
\[
F_{\mathrm{eff}}^2 = a_0^2 + \frac{1}{2}\sum_{\substack{n=1\\n\,\text{impair}}}^{+\infty}b_n^2
= 0 + \frac{1}{2}\sum_{\substack{n=1\\n\,\text{impair}}}^{+\infty}\frac{16}{n^2\pi^2}
= \frac{8}{\pi^2}\sum_{\substack{n=1\\n\,\text{impair}}}^{+\infty}\frac{1}{n^2}
\]
En égalisant les deux expressions :
\[
\frac{8}{\pi^2}\sum_{\substack{n=1\\n\,\text{impair}}}^{+\infty}\frac{1}{n^2} = 1
\qquad\Longrightarrow\qquad
\boxed{\sum_{\substack{n=1\\n\,\text{impair}}}^{+\infty}\frac{1}{n^2} = \frac{\pi^2}{8}}
\]
\end{tcolorbox}

\end{enumerate}

\vfill
\begin{center}
\small\color{gray}\textit{Lycée Denis Diderot, Bavilliers — Alexis RUIZ}
\end{center}

\end{document}
